\section{Método}

En concordancia con el esquema estándar de reporte científico promovido por APA 7, esta sección detalla los procedimientos metodológicos, la caracterización de los participantes o fuentes de datos y las técnicas empleadas.

\subsection{Diseño de la Investigación}
Se empleó un enfoque metodológico mixto, orientado a la sistematización de directrices tipográficas y a la verificación automatizada de consistencia sintáctica en documentos estructurados.

\subsection{Participantes y Muestra}
Para ilustrar la estructura estándar de reporte, se consideró una muestra hipotética de $N = 100$ participantes universitarios divididos equitativamente en dos grupos de evaluación.

\subsection{Instrumentos y Materiales}
Se utilizaron instrumentos de registro estandarizados y matrices de cotejo diseñadas específicamente para medir el apego a las directrices de APA 7.ª edición \parencite{apa2020}.

\subsection{Procedimiento}
El procedimiento se desarrolló a lo largo de tres etapas consecutivas:
\begin{enumerate}
    \item \textbf{Fase de configuración inicial:} Definición de preámbulo, librerías y metadatos del documento.
    \item \textbf{Fase de redacción modular:} Generación de contenidos por sección con citación cruzada.
    \item \textbf{Fase de auditoría editorial:} Escrutinio con la lista de verificación y depuración de referencias bibliográficas.
\end{enumerate}
